%% ============================================================
%% preamble/code.tex — listings environments (module: code)
%%
%% Purpose      : the authoring-contract code surface:
%%                  codelisting  — \begin{codelisting}[language=Python,
%%                                 caption={...}] ... \end{codelisting}
%%                  promptcode   — verbatim human->model prompt
%%                  outputcode   — verbatim program/model output
%%                  \code{...}   — inline code (\lstinline-backed here;
%%                                 commands.tex supplies a \texttt
%%                                 version when this module is off)
%% Dependencies : listings (loaded here), colors.tex (code-* tokens),
%%                fonts.tex (mono face). Loaded only when
%%                \ifBookModuleCode is true (book.yaml modules.code).
%% Provenance   : vibe-coding-for-lawyers' code stack, including the
%%                hard-won \lst@AddToHook catcode fix for `#` inside
%%                listings (docs/research/vibe-coding-for-lawyers.md);
%%                prompt/output variants from the same book's AI
%%                dialogue rendering; 72-char guidance from the RFC
%%                book's CODE-STYLE rules.
%% ============================================================

\usepackage{listings}

% THE `#` FIX (vibe-coding). Inside styles/hooks listings can see `#`
% with its parameter catcode and die with "Illegal parameter number".
% Forcing catcode 12 (other) at listing init makes Python/bash
% comments and shebangs safe everywhere, including inside tcolorbox.
\makeatletter
\lst@AddToHook{Init}{\catcode`\#=12\relax}
\makeatother

% ----------------------------------------------------------------------
% Base style — print-safe by construction:
%   * breaklines with a visible hooked-arrow continuation marker, so
%     a wrapped line can never be misread as two statements;
%   * columns=fullflexible + keepspaces: natural mono spacing without
%     the ragged inter-letter gaps of aligned columns;
%   * WIDTH GUIDANCE: 72 characters is the classic source ceiling, but
%     the 6x9 measure at \small mono fits ~60; keep lines nearer 60 to
%     avoid wraps entirely (wrapped lines are safe, just less pretty).
% ----------------------------------------------------------------------
\lstdefinestyle{bookcode}{
  basicstyle=\ttfamily\small,
  keywordstyle=\color{code-keyword}\bfseries,
  commentstyle=\color{code-comment}\itshape,
  stringstyle=\color{code-string},
  showstringspaces=false,
  columns=fullflexible,
  keepspaces=true,
  tabsize=4,
  breaklines=true,
  breakatwhitespace=true,
  breakindent=1.5em,
  postbreak=\mbox{{\color{code-comment}\footnotesize$\hookrightarrow$}\space},
  frame=leftline,
  framerule=1.5pt,
  rulecolor=\color{code-frame},
  backgroundcolor=\color{code-bg},
  xleftmargin=12pt,
  framexleftmargin=6pt,
  aboveskip=0.9\baselineskip,
  belowskip=0.9\baselineskip,
  captionpos=b,
  upquote=true,               % straight quotes in code, not curly
}

% Captions under listings: quiet, muted, like figure captions.
\makeatletter
\def\lst@makecaption#1#2{% mirror LaTeX's \@makecaption, muted color
  \vskip\abovecaptionskip
  {\small\color{bk-caption}#1: #2\par}%
  \vskip\belowcaptionskip}
\makeatother

% --- codelisting: the one general-purpose listing environment. -------
% Accepts all listings keys; authoring contract blesses language= and
% caption= (label= also works if a book chooses to number listings).
\lstnewenvironment{codelisting}[1][]
  {\lstset{style=bookcode, #1}}
  {}

% --- promptcode: what a human typed at an AI (or shell) prompt. -------
% Accent-tinted, no language highlighting — prompts are prose, and
% keyword-coloring English words looks like a bug in print.
\lstnewenvironment{promptcode}[1][]
  {\lstset{style=bookcode, language={},
           backgroundcolor=\color{prompt-bg},
           rulecolor=\color{prompt-frame}, #1}}
  {}

% --- outputcode: what came back. Neutral tint, thinner rule, ---------
% slightly smaller: output is evidence, not text to study.
\lstnewenvironment{outputcode}[1][]
  {\lstset{style=bookcode, language={},
           basicstyle=\ttfamily\footnotesize,
           backgroundcolor=\color{output-bg},
           rulecolor=\color{output-frame},
           framerule=0.8pt, #1}}
  {}

% --- \code — inline code span. ----------------------------------------
% \lstinline-backed so `#`, `_`, `$` etc. need no escaping in chapters:
% \code{dict['key']} just works. Fragile in moving arguments
% (footnotes, captions) — the style checker flags that usage.
\newcommand{\code}{\lstinline[basicstyle=\ttfamily]}
